\SourcePage{255}{258}
\section{Spinors}

For zero spin the wave function has a single component, $\psi(0)$. The spin
operator annihilates it: $\hat{\mathbf s}\psi=0$. Since $\mathbf s$ is
connected with the operator of infinitesimal rotations, this means that the
wave function of a spin-zero particle is unchanged by rotations of the
coordinate system; it is a scalar.

The wave function of a spin-$1/2$ particle has two components, $\psi(1/2)$
and $\psi(-1/2)$. For convenience in later generalizations, label these by
upper indices 1 and 2. The two-component quantity
\begin{equation}
 \psi=\binom{\psi^1}{\psi^2}=
 \binom{\psi(1/2)}{\psi(-1/2)}
 \tag{56.1}
\end{equation}
is called a \emph{spinor}.

Under an arbitrary rotation of the coordinate system, its components undergo
the linear transformation
\begin{equation}
 \psi^{1'}=a\psi^1+b\psi^2,\qquad
 \psi^{2'}=c\psi^1+d\psi^2.
 \tag{56.2}
\end{equation}
This can be written
\begin{equation}
 \psi^{\lambda'}=(\hat U\psi)^\lambda,\qquad
 \hat U=\begin{pmatrix}a&b\\c&d\end{pmatrix},
 \tag{56.3}
\end{equation}
where $\hat U$ is the transformation matrix\footnote{The notation
$\hat U\psi$ means multiplication of the rows of $\hat U$ by the column
$\psi$.}. Its elements are generally complex functions of the rotation angles
of the coordinate axes. Relations among them follow directly from the
physical requirements on a spinor as a particle wave function.

Consider the bilinear form
\begin{equation}
 \psi^1\varphi^2-\psi^2\varphi^1,
 \tag{56.4}
\end{equation}
where $\psi$ and $\varphi$ are two spinors. Direct calculation gives
\[
 \psi^{1'}\varphi^{2'}-\psi^{2'}\varphi^{1'}
 =(ad-bc)(\psi^1\varphi^2-\psi^2\varphi^1),
\]
so (56.4) transforms into itself under a coordinate rotation. A single
function that transforms into itself can be regarded
\SourcePage{256}{259}as belonging to spin zero; it must consequently be a
scalar and remain altogether unchanged by a
rotation. Hence
\begin{equation}
 ad-bc=1;
 \tag{56.5}
\end{equation}
the determinant of the transformation matrix is unity\footnote{A transformation of two quantities of this kind is called \emph{binary}.
}.


Further relations follow from requiring the expression
\begin{equation}
 \psi^1\psi^{1*}+\psi^2\psi^{2*},
 \tag{56.6}
\end{equation}
which gives the probability of finding the particle at a given point, to be a
scalar. A transformation preserving the sum of squared moduli is unitary, so
$\hat U^+=\hat U^{-1}$ (see \S~12). Subject to (56.5),
\[
 \hat U^{-1}=\begin{pmatrix}d&-b\\-c&a\end{pmatrix}.
\]
Equating this to the conjugate matrix
\[
 \hat U^+=\begin{pmatrix}a^*&c^*\\b^*&d^*\end{pmatrix}
\]
gives
\begin{equation}
 a=d^*,\qquad b=-c^*.
 \tag{56.7}
\end{equation}

By (56.5) and (56.7), the four complex quantities $a,b,c,d$ contain only
three independent real parameters, corresponding to the three angles that
specify a rotation of a three-dimensional coordinate system.

Comparing the scalar expressions (56.4) and (56.6), we see that
$\psi^{1*},\psi^{2*}$ must transform as $\psi^2,-\psi^1$. Relations (56.5)
and (56.7) make this easy to verify\footnote{This property is closely
connected with time-reversal symmetry. Time reversal corresponds (see
\S~18) to replacing a wave function by its complex conjugate, while also
reversing angular-momentum projections. Thus the functions complex conjugate
to $\psi^1=\psi(1/2)$ and $\psi^2=\psi(-1/2)$ must have the same properties
as the components belonging respectively to spin projections $-1/2$ and
$1/2$.}.

Spinor algebra can be put into a form analogous to tensor algebra. Alongside
the contravariant components $\psi^1,\psi^2$, with upper indices, introduce
\SourcePage{257}{260}\emph{covariant} components, with lower indices, by
\begin{equation}
 \psi_1=\psi^2,\qquad \psi_2=-\psi^1.
 \tag{56.8}
\end{equation}
The invariant combination (56.4) of two spinors then becomes the scalar
product
\begin{equation}
 \psi_\lambda\varphi^\lambda
 =\psi_1\varphi^1+\psi_2\varphi^2.
 \tag{56.9}
\end{equation}
Here and below, repeated dummy indices are summed as in tensor algebra. One
rule specific to spinor algebra must be kept in mind. Since
$\psi^\lambda\varphi_\lambda=\psi^1\varphi_1+\psi^2\varphi_2
=-\psi^2\varphi^1-\psi^1\varphi^2$, we have
\begin{equation}
 \psi^\lambda\varphi_\lambda=-\psi_\lambda\varphi^\lambda.
 \tag{56.10}
\end{equation}
It follows at once that the scalar product of any spinor with itself vanishes:
\begin{equation}
 \psi_\lambda\psi^\lambda=0.
 \tag{56.11}
\end{equation}

As explained above, $\psi_1,\psi_2$ transform as
$\psi^{1*},\psi^{2*}$, that is,
\begin{equation}
 \psi'_{\lambda}=(\hat U^*\psi)_\lambda.
 \tag{56.12}
\end{equation}
The product $\hat U^*\psi$ may equally be written $\psi\hat U^*$ using the
transpose of $\hat U^*$. Unitarity gives
$\hat U^*=\widetilde{\hat U}^{-1}$, so that
$\widetilde\psi'=(\widetilde\psi\hat U^{-1})$, or\footnote{Notation of the
form $\psi\hat U$, with $\psi$ to the left, means multiplying the row of
components $(\psi_1,\psi_2)$ by the columns of $\hat U$.}
\begin{equation}
 \psi'_\lambda=\psi_\mu(U^{-1})^\mu{}_{\lambda}.
 \tag{56.13}
\end{equation}

Just as vectors lead to tensors in ordinary tensor algebra, one can introduce
\emph{spinors of higher rank}. A second-rank spinor is a four-component
quantity $\psi^{\lambda\mu}$ whose components transform like the products
$\psi^\lambda\varphi^\mu$ of components of two first-rank spinors. Alongside
the contravariant components $\psi^{\lambda\mu}$ one may consider covariant
$\psi_{\lambda\mu}$ and mixed $\psi_\lambda{}^\mu$ components, transforming
respectively like $\psi_\lambda\varphi_\mu$ and
$\psi_\lambda\varphi^\mu$. Spinors of any rank are defined analogously.

\SourcePage{258}{261}
The passage from contravariant to covariant components and back can be written
\begin{equation}
 \psi_\lambda=g_{\lambda\mu}\psi^\mu,\qquad
 \psi^\lambda=g^{\mu\lambda}\psi_\mu,
 \tag{56.14}
\end{equation}
where
\begin{equation}
 g_{11}=g_{22}=0,\qquad g_{12}=-g_{21}=1
 \tag{56.15}
\end{equation}
is the \emph{metric spinor} in a two-dimensional vector space\footnote{The
matrix (56.15) coincides with $i\hat\sigma_y$.}. Similarly, for example,
\[
 \psi_\lambda{}^\mu=g_{\lambda\nu}\psi^{\nu\mu},\qquad
 \psi_{\lambda\mu}=g_{\lambda\nu}g_{\mu\rho}\psi^{\nu\rho},
\]
so that $\psi_1{}^2=-\psi^{11}$, $\psi_2{}^1=-\psi^{21}$,
$\psi_{11}=\psi^{22}$, and so forth.

The components $g_{\lambda\mu}$ themselves form the antisymmetric unit
spinor of rank two. Their values are unchanged under coordinate
transformations, and one readily verifies
\begin{equation}
 g_{\lambda\nu}g^{\mu\nu}=\delta_\lambda^\mu,
 \tag{56.16}
\end{equation}
where $\delta_1^1=\delta_2^2=1$ and
$\delta_1^2=\delta_2^1=0$.

As in ordinary tensor algebra, spinor algebra has two basic operations:
multiplication and contraction over a pair of indices. Multiplication of two
spinors gives one of higher rank; for example, second- and third-rank spinors
$\psi_{\lambda\mu}$ and $\varphi^{\nu\rho\sigma}$ give the fifth-rank
spinor $\psi_{\lambda\mu}\varphi^{\nu\rho\sigma}$. Contraction over a pair
of indices, meaning summation over equal values of one covariant and one
contravariant index, lowers the rank by two. Thus contraction of
$\psi_{\lambda\mu}{}^{\nu\rho\sigma}$ over $\mu$ and $\nu$ gives the
third-rank spinor $\psi_{\lambda\mu}{}^{\mu\rho\sigma}$; contraction of
$\varphi_\lambda{}^\mu$ gives the scalar
$\varphi_\lambda{}^\lambda$. A rule analogous to (56.10) holds: interchanging
the upper and lower positions of the contracted indices changes the sign,
so $\varphi_\lambda{}^\lambda=-\varphi^\lambda{}_{\lambda}$. In particular,
contracting two indices over which a spinor is symmetric gives zero. Thus for
a symmetric second-rank spinor $\varphi_\lambda{}^\mu$, one has
$\varphi_\lambda{}^\lambda=0$.

A \emph{symmetric spinor} of rank $n$ is symmetric in all its indices. A
nonsymmetric spinor can be symmetrized by summing the components obtained from
all possible permutations of its indices. From the preceding result, no
spinor of lower rank can be formed by contracting components of a symmetric
spinor.

\SourcePage{259}{262}
A spinor antisymmetric in all its indices can only have rank two. Indeed, each
index has only two possible values; with three or more indices at least two
must have equal values, and every component then vanishes identically. Every
antisymmetric second-rank spinor is a scalar multiple of the unit spinor
$g_{\lambda\mu}$. One consequence is
\begin{equation}
 g_{\lambda\mu}\psi_\nu+g_{\mu\nu}\psi_\lambda
 +g_{\nu\lambda}\psi_\mu=0,
 \tag{56.17}
\end{equation}
where $\psi_\lambda$ is arbitrary. The left-hand side is simply, as can be
checked directly, an antisymmetric third-rank spinor and must therefore
vanish.

The spinor formed by multiplying $\psi_{\lambda\mu}$ by itself and
contracting one pair of indices is antisymmetric in the other pair, since
\[
 \psi_{\lambda\nu}\psi_\mu{}^\nu
 =-\psi_\lambda{}^\nu\psi_{\mu\nu}.
\]
It must consequently reduce to $g_{\lambda\mu}$ times a scalar. Determining
that scalar by requiring contraction over the second pair to give the proper
result, we obtain
\begin{equation}
 \psi_{\lambda\nu}\psi_\mu{}^\nu
 =-\frac12\psi_{\rho\sigma}\psi^{\rho\sigma}g_{\lambda\mu}.
 \tag{56.18}
\end{equation}

The components of the spinor $\varphi^{\lambda\mu\ldots}$, complex conjugate to
the spinor $\psi_{\lambda\mu\ldots}$, transform as the components of the
contravariant spinor $\varphi^{\lambda\mu\ldots}$, and conversely.

Consequently, for any spinor the sum of the squared moduli of its components is an invariant.

